\documentclass[11pt]{article}
\usepackage{amsmath,amssymb,amsthm,booktabs,hyperref}
\newtheorem{theorem}{Theorem}
\newtheorem{lemma}[theorem]{Lemma}
\newtheorem{proposition}[theorem]{Proposition}
\newtheorem{corollary}[theorem]{Corollary}
\newcommand{\sq}{\sqcup}
\title{Ordinary gluing for the remaining centre types}
\author{Human Compression project}
\date{}
\begin{document}
\maketitle

\paragraph{Scope and dependencies.}
We consider connected simple nonbipartite quintic graphs of order $32$,
girth four and diameter eight, called target graphs below. The original
cage problem and its published computational census are due to Cambie,
Goedgebeur, Jooken and Van den Eede~\cite{original}. We import the
retained fixed-geodesic reduction~\cite{middle}, the complete twenty-four
nonisolated middle rows~\cite{rows}, the intrinsic centre decomposition
and gluing theorem~\cite{centre}, and the three centre strata proved
in~\cite{halves}. Their hypotheses and ordinary isomorphism qualifications
are retained.

The present contribution is the ordinary identification of the other
centre components, their reuse as common half domains, and their product
formulas. The constructions were suggested in the project discussion;
their proofs are written here in full. Standard graph isomorphism,
independent subsets, integer margins and graph distance remain standard
ingredients. No priority over the original census is claimed. No subject
Source, graph enumeration, TeX compilation or Lean verification was run
in preparing this note. Exact small orbit counts below are analytic
theorems, not computational validation receipts.

\section{Notation and gluing}

The centre has one or two full-graph false-twin classes, and deleting it
leaves two connected components of order at most sixteen. An outside
vertex meets all or none of each centre class. A boundary color specifies
the classes it meets. If a boundary vertex meets centre weight $w$, its
internal degree is $5-w$. In a joined two-class centre, no outside vertex
can meet both classes, since this would create a triangle.

When the two centre classes have distinct sizes we order them by size.
For equal sizes, a class interchange is allowed only simultaneously in
the two components. The intrinsic gluing theorem says that these weighted
centre data and the pair of boundary-colored ordinary components give a
complete ordinary uncolored graph invariant. In particular there are
no separate matchings of ports to individual vertices of a twin class.

We use the retained fixed-geodesic notation only to prove coverage and
converses: the independent neighborhoods $A,B,C$ have size five each;
their complementary residual graph has five edges; each endpoint has
five row neighbors, and every row has four side neighbors. The endpoint
columns of each original side are positive. Each original side has at
least five vertices. The middle count $q$ and the row labels $Q_i$ are
not ordinary invariants.

The half domains $\mathcal H_b$ ($1\leq b\leq4$) and counts $h_b$ are
those of~\cite{halves}: an endpoint $p$, five rows, a six-vertex side
$D$ with $5-b$ edges, and $b$ intrinsic degree-two ports. Rows are
independent four-subsets, ports are independent two-subsets, and
\[
 \alpha(v)+\beta(v)+d_D(v)=5,\qquad \alpha(v)\geq0.
\]
The ports have no other edges. These are ordinary unrooted half counts,
not numbers of endpoint-marked records. All other twelve vertices have
internal degree five. A zero column, if present, is unique; $D$ is then
a star with that vertex as centre and every port contains it. The complete
side types and a direct radius-four converse are already proved there.

Finally let $H$ denote the ordinary count of the enlarged sixteen-vertex
three-port domain in~\cite{centre}: its side has seven vertices and three
edges, its five rows are independent four-subsets, its three ports are
independent three-subsets, and its endpoint columns may be zero. The
complete one-class two-centre stratum has count $6H$. This symbol does
not assert a numerical evaluation of the enlarged domain.

\section{Two small half counts}

\begin{lemma}\label{zeroinvariant}
An ordinary half in a domain $\mathcal H_b$ cannot have both a positive
endpoint-column presentation and a presentation with a zero column.
\end{lemma}
\begin{proof}
Use a presentation with zero column $u$. The side is a star centred at
$u$, with $e=5-b$ leaves and $b$ side isolates. Every port meets $u$
and one side isolate. The five other side vertices form an independent
set $S$; every endpoint row meets four of them. All other side columns
are positive.

Consider any candidate endpoint $p'$ in another presentation. The ports
are intrinsic degree-two vertices. An endpoint has degree five, does not
meet a port, and none of its five row neighbors meets a port, since rows
and ports have no edges between them. A port-supported side vertex,
including $u$, therefore cannot be $p'$. An unported star leaf cannot
be $p'$ either, since its neighborhood contains the port-supported $u$.
An unported side isolate meets all five old rows, so is a false twin
of the old endpoint $p$; choosing it in place of $p$ leaves the zero
column at $u$. Choosing $p$ itself also leaves that zero column.

It remains to choose an old row $a$. Its neighborhood is $\{p\}$ and
four vertices of $S$. Let $t$ be the omitted member of $S$. In the new
presentation the rows must be exactly that neighborhood. The vertex
$t$ meets neither $p$ nor another member of $S$, since $S$ is independent.
Thus $t$ has no new endpoint-row incidence. It is neither the chosen
endpoint, one of its neighbors nor a port, so is a new side vertex
with zero column. We have exhausted every degree-five vertex. No
candidate can give a positive presentation, proving the assertion.
\end{proof}

\begin{proposition}\label{hsmall}
The ordinary half counts satisfy $h_1=2$ and $h_2=5$.
\end{proposition}
\begin{proof}
For $b=1$ the side types are $C_4\sq2K_1$ and $K_{1,4}\sq K_1$.
In the cycle type, every endpoint four-subset contains both isolates
and an opposite cycle pair. The isolates have no port incidence. The
single independent two-subset port is one opposite cycle pair; the
five row pairs then have multiplicities two for that pair and three
for the other, by the degree equations. This gives exactly one
endpoint-marked type, hence at most one ordinary half. In the star
type the centre column is zero and its port must contain the centre
and the sole isolate. Each of the other five vertices is omitted once
among the rows, giving exactly one further type. Both satisfy all
margins. Lemma~\ref{zeroinvariant} separates them ordinarily, so $h_1=2$.

For $b=2$ the side types are $P_4\sq2K_1$ and $K_{1,3}\sq2K_1$.
First label the path vertices $0,1,2,3$ in order. Every independent
four-subset of the side consists of the two isolates and one of the
three path pairs
\[
 X=\{0,2\},\qquad Y=\{0,3\},\qquad Z=\{1,3\}.
\]
The side isolates meet all five rows, so no port can use them. Thus
each port uses one of $X,Y,Z$ as well. Let $m_X,m_Y,m_Z$ be the combined
multiplicities among the five row pairs and two port pairs. The side
degrees require
\[
 m_X=3,\qquad m_Z=3,\qquad m_X+m_Y=4,
\]
and hence $(m_X,m_Y,m_Z)=(3,1,3)$. Up to reversal of the path, which
interchanges $X,Z$, the unordered pair of port choices is exactly
\[
 \{X,X\},\qquad \{X,Z\},\qquad \{X,Y\}.
\]
Two $Y$ choices are impossible since its total multiplicity is one.
In each of the three displayed cases, subtracting the port multiplicities
from $(3,1,3)$ gives valid nonnegative row multiplicities, and every side
column stays positive. Equal choices give equivalent marked records by
permuting the equal-role rows and ports. Distinct choices remain distinct
ordinarily: the two intrinsic port neighborhoods intersect in respectively
two, zero and one vertices. This gives exactly three positive half types.

In the star type the centre $u$ has zero column. Both ports have the
form $\{u,t\}$ with $t$ one of the two isolates. They may use the same
isolate or different isolates, up to the star and isolate symmetries.
The five endpoint rows are four-subsets of the five vertices other than
$u$. Each leaf is omitted once; each isolate is omitted with multiplicity
equal to its port incidence. These prescriptions sum to five and satisfy
all margins, so both choices occur. Their port-neighborhood intersections
have sizes two and one respectively, distinguishing them intrinsically.
Lemma~\ref{zeroinvariant} separates both zero types from every positive
type. Thus there are exactly $3+2=5$ ordinary halves.
\end{proof}

It follows from the one-class centre-three theorem of~\cite{halves} that
this stratum has count
\begin{equation}\label{oneclass3}
 2h_4+5h_3.
\end{equation}
This substitution evaluates neither $h_3$ nor $h_4$.

\section{An independent centre with two twin classes}

\subsection{Centre size three}

The two distinct centre classes have sizes one and two. They are
independent of one another, so their induced centre graph is empty.
The retained centre table restricts this stratum to the middle-edge
$q=3$ rows and the $q=4$ middle-edge rows with an extra middle-to-side
edge. In both cases the components have orders thirteen and sixteen.

\begin{theorem}\label{emptythree}
The ordinary count for this centre stratum is $h_4$. Its thirteen-vertex
component is the unique colored component from the centre-$P_3$ stratum;
its sixteen-vertex component is an arbitrary half in $\mathcal H_4$,
with each of its four ports attached to both centre classes.
\end{theorem}
\begin{proof}
Use the middle-edge notation from~\cite{rows}: a degree-four middle
vertex $u$ has a sole own-side port neighbor $a$, and $r$ neutral middle
leaves, where $r=1$ or $2$. Put $V$ for these leaves and $X$ for the
other centre class. The sizes are $|V|=r$, $|X|=3-r$. The edge $ua$
stays in the thirteen-component. The vertex $u$ meets centre $V$ only
and has $4-r$ side neighbors; $a$ meets centre $X$ only and has $r+1$
side neighbors. Triangle-freeness makes these side supports disjoint.
Their sizes sum to five, and the side has no internal edge. Every side
vertex therefore has four endpoint-row neighbors; the five rows omit
each side vertex once. Ordered by the singleton and double centre
weights, the two exclusive ports have internal degrees four and three,
an edge between them, and complementary side supports of sizes three
and two. This is exactly the unique colored thirteen-component already
proved in~\cite{halves}. The argument is unchanged when $r=1$ and $r=2$
exchange its original middle and port provenance.

In a $q=3$ row the opposite side has size six and one internal edge.
Its four ports each meet all three centre vertices and have two side
neighbors. The twelve nonports retain degree five and the endpoint
columns are positive. This is the positive part of $\mathcal H_4$.
In a $q=4$ row the opposite side has size five and one additional middle
vertex $w$, with one side edge and all four opposite ports as neighbors.
Include $w$ in the six-vertex side. It is a zero endpoint column, the
side has one edge, and each of the four ports has two side neighbors.
This is exactly the zero part of $\mathcal H_4$.

Conversely every positive $\mathcal H_4$ record gives the opposite leg
of a $q=3$ middle-edge row with $r=1$, and every zero record gives the
opposite leg of its $q=4$ extra-middle-edge row. The latter assertion
uses the zero-column characterization: the one side edge meets the
zero vertex, every port meets it, and the other port neighbors lie among
the four side isolates. The unique thirteen-component supplies the
required degree-four middle vertex with three side leaves and its sole
own-side port. All local rows, columns and triangle tests of~\cite{rows}
are then satisfied. That construction and the centre theorem give a
target graph with precisely these two independent centre classes.

The unequal centre weights fix both boundary coordinates, and the
unequal component orders fix their roles. The four common ports of the
sixteen-component all have the same boundary color, which is intrinsic
through their internal degree two. Its ordinary half count is $h_4$.
There is only one colored thirteen-component and no attachment alignment.
The complete gluing theorem therefore gives the stated count.
\end{proof}

\subsection{Centre size four}

Now the independent centre classes have sizes $(1,3)$ or $(2,2)$.
The centre table restricts this case to the $q=4$ middle-edge rows whose
extra residual edge lies inside the opposite side. Both components have
order fourteen.

\begin{lemma}\label{leftemptyfour}
One component has two exclusive ports $u,a$ joined by an edge, one
common port $w$ of degree one, and five independent side vertices $S$.
Their endpoint rows are forced. For a fixed ordered pair of centre
classes there are exactly three colored ordinary types, according as
the unique neighbor of $w$ belongs to $N_S(u)$, to $N_S(a)$, or to
neither.
\end{lemma}
\begin{proof}
Write the neutral middle leaf class as $V$, of size $r\in\{1,2,3\}$,
and the other class as $X$, of size $4-r$. The degree-four middle
vertex $u$ meets $V$ only, its sole original port neighbor $a$ meets
$X$ only, and the other own-side port $w$ meets all four centre vertices.
The edge $ua$ is present. Inside the component, $u$ has $4-r$ side
neighbors, $a$ has $r$ side neighbors, and $w$ has one. The two supports
$U=N_S(u)$ and $W=N_S(a)$ are disjoint by triangle-freeness and cover
four of the five independent side vertices. There is a unique unused
side vertex $t$.

The neighbor of $w$ may lie in $U$, in $W$, or be $t$. All three
categories are possible since both support sizes are positive. The total
side incidence of these three ports is five, with multiplicities at most
two. Thus the five endpoint four-subsets are forced: a side vertex is
omitted as often as its port incidence. They satisfy positive columns
and every row is independent.

The exclusive boundary colors identify $u,a$, and the common boundary
color identifies $w$. Whether the unique neighbor of $w$ meets $u$, meets
$a$, or meets neither is an intrinsic colored graph invariant, regardless
of the selected endpoint. Conversely equal categories allow a permutation
of $S$ preserving $U,W$ and the neighbor of $w$; matching rows by their
omitted vertex extends it to a component isomorphism. These three types
are therefore complete and pairwise distinct as colored ordinary halves.
\end{proof}

\begin{theorem}\label{emptyfour}
The centre with two independent false-twin classes has six ordinary
graphs for class sizes $(1,3)$ and four for class sizes $(2,2)$.
\end{theorem}
\begin{proof}
The other fourteen-component has three common ports of internal degree
one, a five-vertex side with one internal edge, and five endpoint rows.
It is exactly one of the two three-port components in~\cite{halves}.
All three boundary vertices meet both centre classes, so those two
ordinary component types have no further boundary-color distinction.
The exclusive-port component of Lemma~\ref{leftemptyfour} cannot be
exchanged with this common-port component, despite their equal orders.

For centre weights one and three the boundary coordinates are intrinsic,
and the three exclusive-port types remain distinct. All $3\cdot2$ pairs
occur in the imported $r=1$ row, since its independent side supports and
the arbitrary common-port singleton give exactly the three choices,
while the opposite edge gives both port splits. The complete construction
and centre theorem supply target graphs with the prescribed centre.
Ordinary isomorphisms preserve the two colored halves, and component
isomorphisms glue. There are exactly six ordinary graphs.

For equal weights two, interchanging centre classes exchanges $u,a$ in
the exclusive-port component. It exchanges its first two categories
and fixes the unused-side category. Thus these three types give exactly
two orbits under the allowed simultaneous class swap. On the other
component every boundary color is common and is unchanged by that swap.
Its two intrinsic port-split types remain distinct. Every pair occurs
in the $r=2$ row, and the gluing theorem leaves no other identification.
The count is consequently $2\cdot2=4$.
\end{proof}

\section{A centre inducing $K_{2,2}$}

There are two centre classes of size two, with all four cross edges.
This is precisely row $Q_5$ in the retained table. The two components
have order fourteen. Every boundary port meets one centre class only
and has internal degree three. Each component has three such ports:
one of one boundary color and two of the other.

\begin{proposition}\label{k22legs}
For a fixed choice of the singleton port color there are exactly three
colored ordinary component types. Their ports induce $P_3$; the single
port of one color is its middle, and the two other-color ports are its
leaves. If their outside-port neighbor sets are $\{x\},U,W$, then
$x\notin U\cup W$, $|U|=|W|=2$, and the type is determined by
$t=|U\cap W|\in\{0,1,2\}$.
\end{proposition}
\begin{proof}
The $Q_5$ incidence table prescribes a five-vertex independent side,
an endpoint and five rows, together with these three ports inducing
the path. Internal degree three leaves one side neighbor for the
middle port and two for each leaf port. A leaf cannot use $x$, since
the port edge would give a triangle. The two leaf ports are not adjacent,
so their two-subsets of the other four side vertices may intersect in
zero, one or two vertices.

All three intersections occur. The total side port incidence is five,
and no side vertex has more than two such incidences. Row omissions
are forced by these incidences and have total size five, giving positive
side columns and independent rows. Two unordered two-subsets of a
four-set are carried to another such pair by a permutation of the
four-set exactly when their intersection sizes agree. The matching row
omissions then extend that permutation to a component isomorphism.

Conversely the port set and its colors are intrinsic, the middle of
their induced path is intrinsic, and each leaf's two neighbors outside
the port graph are intrinsic. The integer $t$ is therefore preserved
by any colored ordinary isomorphism. This proves both completeness and
nonidentification of the three types.
\end{proof}

\begin{theorem}\label{k22count}
There are exactly six ordinary target graphs with centre $K_{2,2}$.
\end{theorem}
\begin{proof}
Each centre vertex already has two centre neighbors and requires three
outside neighbors. One component has one boundary port of the first
color and two of the second, while the other has the reverse pattern.
All pairs of the three component types are realized by $Q_5$: choose
their two side-subset intersections independently and use the forced
omissions. The full construction theorem and its centre theorem apply.

For ordered centre classes, the two component orientations are fixed.
Interchanging the equal centre classes reverses these orientations;
the two equal-order components can then be exchanged. Hence the ordinary
invariant is the unordered pair of their parameters $t_1,t_2$.
It is complete by the gluing theorem, and it is necessary because the
two intrinsic component parameters must be preserved as a multiset.
There are $\binom{3+1}{2}=6$ such unordered pairs.
\end{proof}

\section{A centre inducing $K_{2,3}$}

Write $C_A$ for its class of size two and $C_B$ for the class of size
three. They are intrinsically ordered by size. An $A$-port meets $C_A$
only and has internal degree three; a $B$-port meets $C_B$ only and has
internal degree two. The centre table gives component orders thirteen
and fourteen. The possible fixed-geodesic rows are $Q_3,Q_6$ and the
internal side edge rows $j=2,3$.

\begin{proposition}\label{k23thirteen}
There are exactly three colored ordinary thirteen-components. Each has
one $A$-port $a$, one $B$-port $b$, and no edge between them. Their
side neighborhoods have sizes three and two; their intersection size
$t\in\{0,1,2\}$ is a complete intrinsic invariant.
\end{proposition}
\begin{proof}
All four imported rows give an endpoint, five rows and a positive
five-vertex side, with these two ports. Here is also a direct elimination
of a possible extra port edge. If the two ports have $k$ edges between
them, then $k=0$ or $1$. Their side arities total $5-2k$, and the
side degree equation gives $2e=25-20-(5-2k)=2k$, so the side has $e=k$
edges. If $k=1$, every endpoint four-subset must omit an endpoint of
the side edge. The other three side vertices meet all five rows and
cannot meet a port. The $A$-port would consequently need two independent
side neighbors among the endpoints of that edge, which is impossible.
Thus $k=e=0$ and the side is independent.

The two port neighborhoods are arbitrary three- and two-subsets of a
five-set. They are not adjacent, so no triangle condition forbids overlap.
The three intersection sizes zero, one and two are all possible and
classify their pairs under the symmetric group of the side. Their
combined incidences sum to five and are at most two per vertex. The
forced omissions give all endpoint rows, with positive side columns.
The distinct internal port degrees identify their roles intrinsically;
their common outside-port neighbors recover $t$. Equal values give an
isomorphism by matching supports and then omitted-vertex rows. Hence
there are exactly three colored ordinary components.
\end{proof}

\begin{proposition}\label{k23fourteen}
There are exactly six colored ordinary fourteen-components. Five have
an induced port graph consisting of one edge and one isolated port;
the sixth has port graph $P_3$.
\end{proposition}
\begin{proof}
The imported rows give an endpoint, five rows, a positive five-vertex
side, one $A$-port $a$ and two $B$-ports $b_1,b_2$. Same-color ports
cannot be adjacent, since they share centre neighbors. Thus their
induced port graph has $k\in\{0,1,2\}$ edges, all from $a$ to a $B$-port.
The total side arity is $7-2k$, and degree balance gives
\[
 2e=25-20-(7-2k)=2k-2.
\]
Hence $k=1$ or $2$, with $e=k-1$.

For $k=1$, name the port edge $ab_1$. The side is independent. Let
$U=N_S(a)$, let $N_S(b_1)=\{x\}$, and let $W=N_S(b_2)$. Then
$|U|=|W|=2$ and $x\notin U$, the latter by triangle avoidance on
$ab_1$. Up to a side permutation the entire choice is determined by
\[
 i=|U\cap W|,\qquad j=\mathbf1_{x\in W}.
\]
There are exactly five possibilities:
\[
 (i,j)=(0,0),(1,0),(2,0),(0,1),(1,1).
\]
Indeed the side partitions as $U$, $\{x\}$, and the remaining two-set;
$W$ takes $i,j,2-i-j$ vertices from these parts. These quantities
determine its orbit, and are feasible exactly for the displayed pairs.
Their total port side incidence is five, with multiplicities at most
two. The forced row omissions therefore satisfy all degree equations
with positive side columns.

The ports are intrinsic by their internal degrees and boundary colors.
Among $b_1,b_2$, the port edge intrinsically distinguishes $b_1$.
Thus $U,\{x\},W$ as outside-port neighbor sets are intrinsic, and the
five pairs $(i,j)$ cannot merge under an unrooted colored isomorphism.
Conversely equal pairs permit a side permutation matching these sets,
which extends over the forced rows. There are exactly five such types.

For $k=2$, the ports induce the star with centre $a$, and the side has
one edge $xy$. Every endpoint row omits $x$ or $y$. Its other three
side vertices have row degree five and cannot meet ports. The three
side arities of $a,b_1,b_2$ are all one, so their neighbors lie among
$x,y$. Triangle avoidance on the port edges forces $a$ to use one
endpoint and both $B$-ports to use the other. Up to interchanging $x,y$
this is unique. Row-omission multiplicities at these endpoints are
$1+1=2$ and $1+2=3$, so the form is feasible and has positive columns.
Its port graph distinguishes it intrinsically from all five previous
types. This proves the total of six ordinary colored components.
\end{proof}

\begin{theorem}\label{k23count}
There are exactly eighteen ordinary target graphs with centre $K_{2,3}$.
\end{theorem}
\begin{proof}
Necessity is given by Propositions~\ref{k23thirteen} and
\ref{k23fourteen}. For sufficiency, all three thirteen-component types
can be paired with each of the five one-edge fourteen-types in $Q_3$.
Use its degree-three middle vertex and its opposite unjoined port as
the $A$- and $B$-ports of the thirteen side. On the other side the
degree-two middle vertex is the $A$-port and chooses the single adjacent
$B$-port; the remaining $B$-port is isolated in the port graph. Their
side arities and the triangle test are exactly $U,\{x\},W$ above.
All three thirteen-types also pair with the port-star fourteen-type
in the $j=2$ internal side edge row: the middle vertex on the edge
side is its $A$-port and the two side ports are its $B$-ports. The
opposite side is the thirteen-component. The forced row omissions
are positive in every case and all stated local triangle tests hold.
Thus the complete construction and intrinsic centre theorems supply
every displayed pair as a target graph with exactly this centre.

The class weights two and three prevent a boundary-coordinate swap,
and the different component orders prevent a component swap. Each
component parameter is intrinsic and complete. Conversely any two
colored component isomorphisms glue over the centre without alignment
choices. The ordinary count is $3\cdot6=18$.
\end{proof}

\section{The complete nonbipartite reduction}

\begin{theorem}\label{total}
Under the imported complete fixed-geodesic and intrinsic centre theorems,
the ordinary count $N_{\rm nb}$ of target graphs is
\begin{equation}\label{totalcount}
 N_{\rm nb}=6H+5h_3+3h_4+46.
\end{equation}
Here $H,h_3,h_4$ are ordinary unrooted counts of the three explicitly
defined nonnegative-column half domains. Formula~\eqref{totalcount}
is a structural reduction, not their numerical evaluation.
\end{theorem}
\begin{proof}
The full centre list, including the retained all-isolated strata, is
exhausted by the following table. Entries are ordinary counts, with
products taken only after the complete colored component identification
proved in the cited statements.
\[
\begin{array}{l|l|l}
 \text{intrinsic centre}&\text{ordinary count}&\text{source}\\\hline
 \text{one twin class of size }2&6H&\text{\cite{centre}}\\
 \text{one twin class of size }3&2h_4+5h_3&\text{\cite{halves}, Proposition~\ref{hsmall}}\\
 \text{one twin class of size }4&6&\text{\cite{halves}}\\
 \text{two independent classes of sizes }1,2&h_4&\text{Theorem~\ref{emptythree}}\\
 \text{two independent classes of sizes }1,3&6&\text{Theorem~\ref{emptyfour}}\\
 \text{two independent classes of sizes }2,2&4&\text{Theorem~\ref{emptyfour}}\\
 P_3&6&\text{\cite{halves}}\\
 K_{2,2}&6&\text{Theorem~\ref{k22count}}\\
 K_{2,3}&18&\text{Theorem~\ref{k23count}}
\end{array}
\]
There is no target with empty middle. A singleton middle has the
two-vertex centre. A larger all-isolated middle has the single centre
class $\{p_4\}\sq M$; its size is three or four. If four middle
vertices were all isolated, every midpoint-neighbor would already have
its five neighbors in $\{p_4\}\sq M$ and could not meet a side,
contradicting the fixed geodesic. The nonisolated cases and their centres
are exactly the imported complete twenty-four-row table. Thus the
displayed intrinsic centre types exhaust all targets.

These nine strata are disjoint as ordinary graph classes. The centre
graph is intrinsic; where its induced graph is independent, its full
open-neighborhood equivalence classes distinguish one twin class from
two and distinguish their weight patterns. Ordinary graph isomorphisms
cannot change these data. Hence their counts add without duplicating
different midpoint or middle-set descriptions. Summing the table gives
\[
 6H+(2h_4+5h_3)+6+h_4+6+4+6+6+18
 =6H+5h_3+3h_4+46,
\]
as required.
\end{proof}

\paragraph{What is and is not counted.}
The constants $2,5,6,4,6,18$ proved here are finite analytic orbit counts.
They do not result from a machine search or from fitting a published
total. The three remaining symbols require complete ordinary unrooted
identification over their own finite domains. In particular $H$ includes
the exceptional zero-column forms; substituting the earlier strict
singleton count would omit part of the two-centre stratum. The code
stream of the strict domain is neither an input nor a stopping rule for
the enlarged one.

An actual computational validation of~\eqref{totalcount} requires reviewed
Source, complete finite-domain receipts, and independent checks of all
ordinary identifications or a complete independent ordinary union. No
such run or Lean certificate is asserted in this written packet, and
no new conjecture, completed manuscript or canonical adoption is claimed.

\begin{thebibliography}{9}
\bibitem{original} Cambie, Goedgebeur, Jooken and Van den Eede,
\emph{On the order-diameter ratio of girth-diameter cages},
arXiv:2511.21144v3, Appendix B, Table 6.
\bibitem{middle} Human Compression project,
\emph{The empty and singleton middle cases for quintic graphs},
SHA256
\texttt{76a5c09c2e3f52536c022572bd95ba611392ad38527c0f2ee40c1e886e9d9eb2},
with its separate independent written review.
\bibitem{rows} Human Compression project,
\emph{Nonisolated middle vertices in quintic graphs of diameter eight},
SHA256
\texttt{522af65bc2bb25d0d09ad68b1c047f2806832230e8452d388ead3483d84f6384},
with separate independent written review.
\bibitem{centre} Human Compression project,
\emph{The intrinsic centre and finite gluing interfaces}, SHA256
\texttt{25987bb9eb742cbc1b8e779408dec9823a99e3a4704e3eadb4f3d6735aa73a3d},
with separate independent written review.
\bibitem{halves} Human Compression project,
\emph{Small halves and three intrinsic centre strata}, SHA256
\texttt{697f8edc67afeeae3f469ff3f4080b11ea7360ce834071d57b75154716699844}.
The documentary manifest specifies exact imported proof and review pins.
\end{thebibliography}
\end{document}
